\documentclass[14pt]{extarticle}
\usepackage[a4paper,margin=1in]{geometry}
\usepackage{times}
\usepackage{setspace}
\setstretch{1.05}
\pagestyle{empty}

\begin{document}

\begin{center}
\textbf{A List of Changes (Significant Improvement on Content)}\\
\textbf{Last updated date: 15 June 2026}
\end{center}

\vspace{0.5em}

\noindent
\textbf{The most significant change}

The evaluation methodology was substantially strengthened by introducing a 
comprehensive statistical assessment framework, including MAE, RMSE, median error, 
MAD, standard deviation, and 95\% bootstrap confidence intervals based on 10,000 resamples. 
This provides a more rigorous and statistically reliable evaluation of the proposed IPM method.

\vspace{0.8em}

\noindent
\textbf{The second most significant change}

Statistical validation of the dynamic evaluation results was added through the 
Shapiro–Wilk normality test and Q–Q plot analysis. The findings show that the 
error distribution deviates from normality, supporting the use of median-based 
metrics and improving the statistical validity of the reported performance analysis.

\vspace{0.8em}

\noindent
\textbf{The third most significant change}

A numerical simulation of the IPM projection model was introduced to analyze the 
geometric relationship between image coordinates and estimated field distance. 
The simulation demonstrates the non-linear depth sensitivity of IPM near the horizon 
and provides a theoretical explanation for the increased estimation error and variance 
observed at long distances.

\vspace{1.2em}

\noindent
\textbf{Deletion if any}

The comparative MAE, RMSE, and related error metrics were removed from the regression-versus-IPM 
evaluation because they could produce misleading conclusions when comparing methods with substantially 
different operating ranges. The comparison was revised to emphasize the maximum reliable estimation 
distance, which more accurately reflects the practical advantage of the proposed IPM approach.

\end{document}
